% 数学ノート集・共通テンプレート v2（積分論2026年度第1部の版面を参考）
% A5判。章内で定義・定理・例・注意等の番号を共通化。
% コンパイル： platex -interaction=nonstopmode math_note_template_v2.tex (2回)
%              dvipdfmx math_note_template_v2.dvi
\documentclass[a5paper,10pt,onecolumn]{jsarticle}
\usepackage[dvipdfmx,top=22mm,bottom=20mm,inner=19mm,outer=18mm,headheight=13pt,headsep=6mm]{geometry}
\usepackage{amsmath,amssymb,amsthm,mathtools}
\usepackage{enumitem,fancyhdr,titlesec}
\usepackage[dvipdfmx,bookmarks=true,colorlinks=false]{hyperref}
\usepackage{pxjahyper}

% ----- 差し替える文書情報 -----
\newcommand{\NoteTitle}{Hopf–Rinowの定理}
\newcommand{\NoteSeries}{数学ノート集}

% ----- 文字・組版 -----
% 参考PDFの質素な明朝＋Computer Modernの印象を引き継ぎながら，
% A5判で窮屈にならないよう本文の行間を少し確保する．
\linespread{1.13}
\setlength{\parindent}{1zw}
\setlength{\parskip}{0pt}
\setlength{\abovedisplayskip}{6pt plus 2pt minus 2pt}
\setlength{\belowdisplayskip}{6pt plus 2pt minus 2pt}
\setlength{\abovedisplayshortskip}{3pt plus 2pt minus 1pt}
\setlength{\belowdisplayshortskip}{4pt plus 2pt minus 1pt}
\setlist[enumerate]{topsep=3pt,itemsep=1pt,parsep=0pt,leftmargin=2.4zw}
\titleformat{\section}{\normalfont\normalsize\bfseries}{\thesection}{0.8zw}{}
\titlespacing*{\section}{0pt}{2.0ex plus .4ex}{.9ex}
\titleformat{\subsection}{\normalfont\normalsize\bfseries}{\thesubsection}{0.7zw}{}
\titlespacing*{\subsection}{0pt}{1.5ex plus .3ex}{.6ex}
\raggedbottom

% ----- 柱とページ番号：ページ上部、罫線なし -----
\fancypagestyle{mathnote}{%
  \fancyhf{}%
  \fancyhead[L]{\footnotesize \NoteSeries\ （\NoteTitle）}%
  \fancyhead[R]{\footnotesize \thepage}%
  \renewcommand{\headrulewidth}{0pt}%
  \renewcommand{\footrulewidth}{0pt}%
}
\pagestyle{mathnote}

% ----- 節内通し番号の定義・定理群 -----
\newtheoremstyle{mathnotetheorem}
 {4.5pt}{4.5pt}{\normalfont}{}%
 {\bfseries}{}{0.65zw}%
 {\thmname{#1}\thmnumber{ #2}\thmnote{（#3）}}
\theoremstyle{mathnotetheorem}
\newtheorem{theorem}{定理}[section]
\newtheorem{definition}[theorem]{定義}
\newtheorem{lemma}[theorem]{補題}
\newtheorem{proposition}[theorem]{命題}
\newtheorem{corollary}[theorem]{系}
\newtheorem{example}[theorem]{例}
\newtheorem{remark}[theorem]{注意}
\renewcommand{\proofname}{証明}
\renewcommand{\qedsymbol}{$\square$}
\makeatletter
\renewenvironment{proof}[1][\proofname]{%
  \par\pushQED{\qed}\normalfont\topsep4\p@\@plus2\p@\relax
  \trivlist\item[\hskip\labelsep\bfseries #1]\ignorespaces
}{\popQED\endtrivlist\@endpefalse}
\makeatother

% ----- 使うノートで適宜拡張する記号 -----
\newcommand{\R}{\mathbb{R}}
\newcommand{\norm}[1]{\lVert #1\rVert}
\newcommand{\abs}[1]{\lvert #1\rvert}

\begin{document}
\section*{読了後にできること}
本稿を読むための前提は，Riemann距離の定義，距離完備性，測地線．
本稿の具体的な到達課題は，具体的なRiemann多様体にHopf--Rinowを適用できるか判定する．


\section{Riemann距離と完備性}
$(M,g)$ を連結な有限次元Riemann多様体とする．
\begin{definition}
二点 $p,q$ のRiemann距離を
\[
 d_g(p,q)=\inf\{L_g(\gamma):\gamma
 \text{ は }p\text{ と }q\text{ を結ぶ区分的滑らかな曲線}\}
\]
で定める．
\end{definition}
\begin{lemma}[Riemann距離の位相]
$d_g$ は距離であり，もとの多様体位相を誘導する．
\end{lemma}
\begin{proof}
連結な滑らかな多様体は道連結であるから，
二点を結ぶ区分的滑らかな曲線が存在し距離は有限である．
非負性・対称性・三角不等式は曲線長の性質から従う．
座標近傍の相対コンパクトな小球上で
$g$ とEuclid計量は定数倍で上下から比較できる．
従って近い二点について
$d_g(p,q)\le C|\varphi(p)-\varphi(q)|$．
一方，$p$ の小座標球から出る任意の曲線は
下から $\sqrt c$ 倍のEuclid半径以上の長さを持つ．
よって $q\ne p$ なら $d_g(p,q)>0$ であり，
$d_g$-小球と座標近傍は互いに包含し合う．
これにより距離性と位相の一致が従う．
\end{proof}
\section{測地線完備性}
\begin{definition}
全ての初期データ $(p,v)\in TM$ の測地線が
時刻全体 $\mathbb R$ へ延長できるとき
$(M,g)$ を測地的完備という．
\end{definition}
\begin{proposition}[距離完備性から測地的完備性]
$(M,d_g)$ が完備なら測地的完備である．
\end{proposition}
\begin{proof}
最大測地線が有限時刻 $b$ で延長できないと仮定する．
測地線の速さ $c$ は一定なので
$t,s\uparrow b$ に対し
$d_g(\gamma(t),\gamma(s))\le c|t-s|$．
従って $\gamma(t)$ はCauchyであり，距離完備性から
$p\in M$ に収束する．
$g$ とEuclid計量の局所同値性および一定速度から，
$p$ の相対コンパクト座標近傍で
$(\gamma(t),\dot\gamma(t))$ は接束のコンパクト部分集合内にある．
測地線方程式の局所存在定理はこのコンパクト集合上で
共通の正の時間幅の解を与えるので，
$b$ 直前から解を延長して $b$ を越えられる．
極大性に矛盾する．負方向も同様である．
\end{proof}
\section{完備性から最短測地線へ}
\begin{lemma}[最短曲線の角は消える]\label{lem:corner}
二本の滑らかな測地線を接続した区分的滑らかな曲線が
その端点間の距離を実現するなら，
接続点における片側速度は一致する
（単位速度でパラメータ付けした場合）．
\end{lemma}
\begin{proof}
接続点を $x$ とし，接続の前後の単位速度を
$u,v\in T_xM$ とする．$u\ne v$ なら
小さい $\varepsilon>0$ について
前後の点はそれぞれ
$\exp_x(-\varepsilon u)$，
$\exp_x(\varepsilon v)$ となる．
正規座標でこれらの二点間の距離は
$\varepsilon|u+v|+o(\varepsilon)$ であり，
$|u+v|<2$ なので $2\varepsilon$ より短い曲線で結べる．
元の曲線の長さ $2\varepsilon$ の部分をその曲線で
置き換えると全長が短くなり最短性に矛盾する．
ゆえに $u=v$ である．
\end{proof}
\begin{theorem}[測地的完備なら最短測地線が存在]
\label{thm:reach}
$(M,g)$ が測地的完備なら任意の二点 $p,q$ は
長さ $d_g(p,q)$ の測地線で結ばれる．
\end{theorem}
\begin{proof}
$q\ne p$ とし $d=d_g(p,q)>0$ と置く．
$p$ の正規球の半径より小さい $0<r<d$ を取る．
半径 $r$ の放射測地線の端点全体
$S_r=\exp_p\{v:|v|=r\}$ はコンパクトである．
任意の $p$ から $q$ への曲線は最初に $S_r$ と交わるため
\[
 d=\inf_{x\in S_r}\{r+d_g(x,q)\}.
\]
$S_r$ のコンパクト性で最小点 $x_0$ を取れば
$d=r+d_g(x_0,q)$ となる．
$p$ から $x_0$ への放射測地線を $\gamma$ とし，
測地的完備性から全時刻へ延長する．

集合
\[
 A=\{t\in[r,d]:d=t+d_g(\gamma(t),q)\}
\]
を考える．$r\in A$ であり，距離の連続性によって $A$ は閉．
$t\in A$，$t<d$ とする．$y=\gamma(t)$ の正規球半径
より小さい $0<\varepsilon<d-t$ を取り，
球面 $S_\varepsilon(y)$ 上で $d_g(z,q)$ を最小にする
$z$ を選ぶ．$y$ から $q$ への任意の曲線はこの球面と交わるので
$d_g(y,q)=\varepsilon+d_g(z,q)$ である．
従って
\[
 d=t+\varepsilon+d_g(z,q).
\]
一方 $d\le d_g(p,z)+d_g(z,q)$ であり，
$p$ から $y$ までの $\gamma$ と
$y$ から $z$ への短い放射測地線を結べば
$d_g(p,z)\le t+\varepsilon$．
従って等号が全て成立し，接続した曲線は
$p$ から $z$ への最短曲線である．
補題\ref{lem:corner}からこの曲線に角はなく，
測地線の一意性により $z=\gamma(t+\varepsilon)$ となる．
よって $t+\varepsilon\in A$．
したがって $A$ は $[r,d]$ の右側に開いている．
$A$ は閉・非空であり，右へ延長可能なので
$\sup A=d$ である．閉性から $d\in A$ となり，
$d_g(\gamma(d),q)=0$．すなわち $\gamma(d)=q$ である．
$\gamma$ は単位速度なので長さは $d$ となる．
\end{proof}
\section{Hopf–Rinowの定理}
\begin{theorem}[Hopf--Rinow]\label{thm:hopf}
連結Riemann多様体 $M$ について次は同値である．
\begin{enumerate}
\item $(M,d_g)$ は完備である．
\item $M$ は測地的完備である．
\item 任意の $p\in M$ について指数写像
 $\exp_p:T_pM\to M$ が全空間で定義される．
\item $M$ の閉有界集合はコンパクトである．
\end{enumerate}
これらが成立すると任意の二点は最短測地線で結ばれる．
\end{theorem}
\begin{proof}
(1)$\Rightarrow$(2) は前命題による．
(2)$\Rightarrow$(3) は測地的完備性の定義による．
(3) なら定理\ref{thm:reach}の証明で用いる測地線の
必要な延長がすべて可能なので，
任意の $q$ は $p$ からの最短測地線によって到達できる．
従って閉距離球
$\overline B_g(p,R)$ は
$\exp_p$ による $T_pM$ の閉球
$\{v:|v|\le R\}$ の像である．
後者は有限次元Euclid空間でコンパクトなので
$\overline B_g(p,R)$ はコンパクトである．
任意の閉有界集合はこのコンパクト球の閉部分集合だから
(4) が従う．
(4) ならCauchy列はある閉距離球に尾部が含まれ，
その球はコンパクトなので収束部分列を持つ．
Cauchy性から列全体がその極限へ収束し (1) を得る．
最短測地線の存在は定理\ref{thm:reach}による．
\end{proof}

\section{到達課題の完全解答}
\noindent\textbf{完備性を距離と測地線から判定する．}
\begin{proposition}[開区間と円周の比較]
標準計量を備えた開区間 $(0,1)$ は距離完備でも
測地的完備でもない．
一方，標準計量を備えた円周 $S^1$ は双方の意味で完備である．
\end{proposition}
\begin{proof}
$(0,1)$ のRiemann距離は $|x-y|$ である．
$x_n=1/n$（$n\ge2$）はCauchy列だが
$(0,1)$ 内に極限を持たない．
また $\gamma(t)=1/2+t$ は測地線であるが
$t=1/2$ で境界へ達し延長できない．
$S^1$ はコンパクトなRiemann多様体なので，
Riemann距離から得る距離空間もコンパクトであり，
特に完備である．Hopf--Rinow定理から
測地的完備性が従う．
\end{proof}
\begin{proposition}[完備な非コンパクトな例]
標準Euclid計量を備えた $\mathbb R^n$ は
完備だがコンパクトではない．
任意の二点 $p,q$ は
$\gamma(t)=(1-t)p+tq$（$0\le t\le1$）という
最短測地線で結ばれる．
\end{proposition}
\begin{proof}
Euclid距離の完備性を用いる．
非コンパクト性は列 $(ke_1)_{k\ge1}$ が
収束部分列を持たないことによる．
任意の区分的滑らかな曲線 $c$ に対し
\[
L(c)=\int|\dot c(t)|\,dt
\ge\left|\int\dot c(t)\,dt\right|=|q-p|
\]
である．表示した線分の長さは $|q-p|$ である．
\end{proof}
\noindent\textbf{到達確認．}
$(0,\infty)$ のEuclid計量は完備か．
\par\noindent\textbf{解答．}
$x_n=1/n$ は領域内のCauchy列だが極限0を含まない．
よって不完備である．
\end{document}
